\clearpage\begingroup\SourceLeaf{359b}\ArithmeticHead{28}{L}{STÄCKEL, GAUSS ALS GEOMETER.}
\renewcommand{\thefootnote}{\arabic{footnote})}\setcounter{footnote}{0}
\begin{center}\textbf{6.}\par\smallskip
\textls[100]{Die Entdeckung der transzendenten Trigonometrie (1805—1817).}
\end{center}
\noindent\hspace*{\parindent}\textsc{Schumacher} ist im Wintersemester 1808/9 in Göttingen gewesen, um sich
bei \textsc{Gauss} zum Astronomen auszubilden; während dieser Zeit hat er Aufzeichnungen
über seine Gespräche mit \textsc{Gauss} gemacht. Diese »\textit{Gaussiana}«
bringen unter dem November 1808 die Bemerkung: »\textsc{Gauss} hat die Theorie
der Parallellinien darauf zurückgebracht, dass wenn die angenommene Theorie
nicht wahr wäre, es eine konstante, a priori der Länge nach gegebene Linie
geben müsste, welches absurd ist. Doch hält er selbst diese Arbeit noch
nicht für hinreichend« (W. VIII, S. 165). Hieraus geht hervor, dass \textsc{Gauss}
auch im Jahre 1808 noch schwankte. »Auf die Worte: ‘welches absurd ist’
wollen wir dabei noch nicht einmal das geringste Gewicht legen, denn es ist
höchst wahrscheinlich, dass \textsc{Schumacher} sie aus seinem Eigenen hinzugefügt
hat, wohl aber legen wir Gewicht auf den nachfolgenden Satz. Wenn \textsc{Gauss}
selber seine Untersuchungen noch nicht für abgeschlossen hielt, so muss er
\clearpage\SourceLeaf{360}\ArithmeticHead{29}{R}{DIE GRUNDLAGEN DER GEOMETRIE.}\setcounter{footnote}{0}
\noindent noch immer halb und halb an \textsc{Euklid} geglaubt haben; auf alle Fälle war er
auch damals noch nicht vollständig von der Unbeweisbarkeit des Parallelenaxioms
überzeugt«\footnote{F. \textsc{Engel}, \textit{Lobatschefskijs Leben und Schriften}, in dem Werke: \textit{N. I. Lobatschefskij, zwei geometrische Abhandlungen}, herausgegeben von F. \textsc{Engel}, Leipzig 1898—99, S. 380; im Folgenden angeführt mit \textsc{Lob.}}.
\par
Dass in der nichteuklidischen Geometrie, in der, ebenso wie in der
Sphärik, die Ähnlichkeit von Figuren aufhört, eine a priori gegebene Einheit
der Länge vorhanden ist, hatte schon 1766 \textsc{Lambert} erkannt (P. Th. S. 200),
und \textsc{Legendre} hatte 1794 auf die angebliche Widersinnigkeit eines solchen
absoluten Masses einen Beweis des Parallelenaxioms gegründet.
\par
Auch eine Bemerkung von \textsc{Gauss} aus dem Jahre 1813 ist wohl in demselben
Sinne aufzufassen: »In der Theorie der Parallellinien sind wir jetzt
noch nicht weiter als \textsc{Euklid} war. Dies ist die partie honteuse der Mathematik,
die früh oder spät eine ganz andere Gestalt bekommen muss« (W. VIII,
S. 166). Man wird dabei an den Ausspruch \textsc{d’Alemberts} vom Jahre 1759
erinnert: »Die Erklärung und die Eigenschaften der geraden Linie sowie der
parallelen Geraden sind die Klippe und sozusagen das Ärgernis (le scandale)
der Elementargeometrie«\footnote{J. \textsc{d’Alembert}, \textit{Mélanges de littérature, d’histoire et de philosophie}, t. V., 4. éd. Amsterdam 1767, S. 200.}.
\par
Ein anderer Ton wird in der Besprechung von zwei Beweisversuchen
angeschlagen, die \textsc{Gauss} in den Göttinger Gelehrten Anzeigen vom 20. April
1816 veröffentlicht hat (W. IV, S. 364, VIII, S. 170). Hier spricht er von
dem »eitelen Bemühen, die Lücke, die man nicht ausfüllen kann, durch ein
unhaltbares Gewebe von Scheinbeweisen zu verbergen«. Dass \textsc{Gauss} damit
auf seine Überzeugung von der Unbeweisbarkeit des elften Axioms hindeuten
wollte, wird durch den S. 8, 9 angeführten Brief an \textsc{Schumacher} vom 15. Januar
1827 bestätigt. Ein weiteres Zeugnis dafür, dass er jetzt zur Gewissheit durchgedrungen
war, ist der Brief an \textsc{Gerling} vom 11. April 1816, also gerade aus
der Zeit, in der er jene Besprechung verfasst hatte (W. VIII, S. 168) \textsc{Gauss}
äussert sich hier, auf \textsc{Gerlings} Wunsch, zu dem vorher erwähnten Beweisversuch
\textsc{Legendres} und sagt: »Es scheint etwas paradox, dass eine konstante
Linie gleichsam a priori möglich sein könne; ich finde aber darin nichts
Widersprechendes. Es wäre sogar wünschenswert, dass die Geometrie \textsc{Euklids}
\clearpage\SourceLeaf{361}\ArithmeticHead{30}{L}{STÄCKEL, GAUSS ALS GEOMETER.}\setcounter{footnote}{0}
\noindent nicht wahr wäre, weil wir dann ein allgemeines Mass a priori hätten, z. B.
könnte man als Raumeinheit die Seite desjenigen gleichseitigen Dreiecks
annehmen, dessen Winkel \(=59^0\,59'\,59''{,}99999\)«.
\par
\textsc{Gauss} durfte sich mit solcher Entschiedenheit äussern, denn er war jetzt
im Besitz der Trigonometrie, die in der nichteuklidischen Geometrie gilt.
Wir wissen dies nicht aus Aufzeichnungen oder Briefen von \textsc{Gauss}, sondern
durch einen glücklichen Zufall. Ebenfalls im April 1816 hatte \textsc{Gauss} den
Besuch seines Schülers \textsc{Wachter} erhalten, der auf der Reise nach Danzig,
wo er Professor am Gymnasium illustre geworden war, Göttingen berührte\footnote{Vgl. hierfür wie für die folgenden Angaben den Aufsatz von P. \textsc{Stäckel}: \textit{F. L. Wachter}, Math. Annalen, Bd. 54, 1901, S. 49—85.}.
\textsc{Wachter} hatte kurz vorher in der Zeitschrift für Astronomie und verwandte
Wissenschaften eine Besprechung derselben Parallelentheorie von \textsc{Metternich}
veröffentlicht, mit der \textsc{Gauss} sich in den Göttinger Nachrichten vom 20. April
1816 beschäftigt hat, und so ist es erklärlich, dass die Unterhaltung sich auch
den Grundlagen der Geometrie zuwandte. Hierauf bezieht sich ein Brief von
\textsc{Wachter} an \textsc{Gauss} vom 12. Dezember 1816, in dem jener über Untersuchungen
berichtet, die er, angeregt durch das Gespräch, über die »antieuklidische
Geometrie« angestellt hatte (W. VIII, S. 175). Wir erfahren
hieraus, dass \textsc{Gauss} ihm von seiner transzendenten Trigonometrie gesprochen
hatte; \textsc{Wachter} hatte sich vergeblich bemüht, einen Eingang in diese zu
finden. Man wird daher annehmen dürfen, dass \textsc{Gauss} damals, wie er in
einem bald darauf, am 16. März 1819, an \textsc{Gerling} geschriebenen Briefe sagt,
die nichteuklidische Geometrie »so weit ausgebildet hatte, dass er alle Aufgaben
vollständig lösen konnte, sobald die Konstante \(=C\) gegeben wird«
(W. VIII, S. 182).
\par
Die jetzt gewonnene feste Stellung gibt sich kund in dem Briefe an
\textsc{Olbers} vom 28. April 1817: »Ich komme immer mehr zu der Überzeugung,
dass die Notwendigkeit unserer Geometrie nicht bewiesen werden kann, wenigstens
nicht \textls[100]{vom menschlichen} Verstande noch für den menschlichen Verstand.
Vielleicht kommen wir in einem andern Leben zu andern Einsichten
in das Wesen des Raums, die uns jetzt unerreichbar sind. Bis dahin müsste
man die Geometrie nicht mit der Arithmetik, die rein a priori steht, sondern
etwa mit der Mechanik in gleichen Rang setzen« (W. VIII, S. 177).
\clearpage\SourceLeaf{362}\ArithmeticHead{31}{R}{DIE GRUNDLAGEN DER GEOMETRIE.}
\noindent\hspace*{\parindent}Auf welchem Wege \textsc{Gauss} zur nichteuklidischen Trigonometrie gelangt
ist, lässt sich nicht mit Sicherheit sagen. Im Nachlass findet sich eine Herleitung
der Formeln, die wahrscheinlich im Jahre 1846 niedergeschrieben ist
(W. VIII, S. 255). In ihr wird das Verfahren der geometrischen Konstruktionen
und daraus hergeleiteten Funktionalgleichungen angewandt, das \textsc{Gauss},
wie wir gesehen haben, schon im Jahre 1803, freilich ohne Erfolg, benutzt
hatte, und es liegt daher nahe anzunehmen, dass er in der Zeit zwischen 1813
und 1816 auf diesem Wege vorgegangen ist; allein man darf hierin nicht
mehr als eine Vermutung erblicken.
\clearpage\endgroup
